An integer $N\ge 1$ is called \emph{representable} if there exist a prime $p$, an integer $a\ge 1$, and an integer $b$ with $0\le b<p$ such that
\[
 N=ap^2+b.
\]
Question: Is every sufficiently large integer representable?

Variant questions (from the problem text):
- What if the condition that $p$ is prime is omitted (i.e. $p$ is any integer $\ge 2$)?
- Define $c_N$ as the minimal constant such that $N=ap^2+b$ with $0\le b<c_N p$ for some prime $p\le\sqrt N$. Is it true that $c_N$ is eventually $\le 1$ (the main question), or perhaps $\limsup c_N=\infty$? Is $c_N<N^{o(1)}$?
